
\documentclass[11pt]{article}
\usepackage[margin=0.85in]{geometry}
\usepackage{xcolor}
\usepackage{hyperref}
\usepackage{microtype}
\usepackage{framed}
\definecolor{beforegray}{RGB}{248,248,248}
\definecolor{nowblue}{RGB}{241,247,255}
\newenvironment{beforebox}{\def\FrameCommand{\fboxsep=7pt\colorbox{beforegray}}\MakeFramed{\advance\hsize-2\width\FrameRestore}\small}{\endMakeFramed}
\newenvironment{nowbox}{\def\FrameCommand{\fboxsep=7pt\colorbox{nowblue}}\MakeFramed{\advance\hsize-2\width\FrameRestore}\small}{\endMakeFramed}
\title{RegTrace Paper V2 Paragraph Change Report}
\author{}
\date{Comparison: previous text at commit af33cc8 vs current V2 at commit f410c72}
\begin{document}
\maketitle
\noindent This report lists only paragraphs whose text changed between the previous paper text and the current V2 text. Formatting-only files and generated artifacts are excluded.


\section*{Changed Paragraph 1}

\noindent\textbf{Before}

\begin{beforebox}
We make four contributions. First, we define \textbackslash{}task\{\} and introduce \textbackslash{}framework\{\}, a closed-loop framework for evidence retrieval, obligation checking, feedback construction, and reviewer optimization. Second, we instantiate it as \textbackslash{}dataset\{\}, a public EDGAR benchmark with SEC comments, company responses, amended-filing evidence, labels, and feedback fields. Third, we show that GEPA with full natural-language feedback outperforms handwritten prompting, MIPROv2, and scalar/category GEPA variants; a guarded verifier distilled from the learned policy improves further. Fourth, we provide evidence ablations, feedback-component ablations, prompt-evolution analysis, and systematic corroboration from real next-round SEC follow-up comments.
\end{beforebox}

\noindent\textbf{Now}

\begin{nowbox}
We make four contributions. First, we define \textbackslash{}task\{\} and introduce \textbackslash{}framework\{\}, a feedback-learning framework for evidence retrieval, obligation checking, feedback construction, and reviewer optimization. Second, we instantiate it as \textbackslash{}dataset\{\}, a public EDGAR benchmark with SEC comments, company responses, amended-filing evidence, labels, and feedback fields. Third, we show that GEPA with full natural-language feedback outperforms handwritten prompting, MIPROv2, scalar/category GEPA variants, and a supervised lexical baseline; a guarded verifier distilled from the learned policy improves further. Fourth, we provide evidence ablations, independent-feedback and feedback-component ablations, prompt-evolution analysis, and systematic corroboration from real next-round SEC follow-up comments.
\end{nowbox}

\section*{Changed Paragraph 2}

\noindent\textbf{Before}

\begin{beforebox}
We compare five methods. Baseline prompt is a handwritten instruction. MIPROv2 is scalar prompt optimization. GEPA-scalar uses reflective prompt optimization with correctness-only feedback. GEPA-category adds coarse category feedback. GEPA-full uses full natural-language adjudication feedback.
\end{beforebox}

\noindent\textbf{Now}

\begin{nowbox}
We compare five prompt methods. Baseline prompt is a handwritten instruction. MIPROv2 is scalar prompt optimization. GEPA-scalar uses reflective prompt optimization with correctness-only feedback. GEPA-category adds coarse category feedback. GEPA-full uses full natural-language adjudication feedback. We also include a TF-IDF logistic-regression baseline and an independent-feedback GEPA ablation for reviewer-facing robustness.
\end{nowbox}

\section*{Changed Paragraph 3}

\noindent\textbf{Before}

\begin{beforebox}
The majority baseline has high accuracy because unresolved examples dominate, but zero resolved F1. The handwritten prompt is conservative: it has perfect unresolved precision but misses most gaps. MIPROv2 improves substantially, GEPA-scalar improves recall, and GEPA-category is too coarse in this run. GEPA-full is strongest overall, improving macro-F1 by 10.4 points over MIPROv2 and 8.1 points over GEPA-scalar.\par
\par
We also compute five-fold out-of-fold results over review-thread groups. Mean macro-F1 is 0.698 for GEPA-full, 0.647 for MIPROv2, 0.609 for GEPA-scalar, 0.545 for GEPA-category, and 0.446 for the handwritten baseline. The fold standard deviations are nontrivial (0.151 for GEPA-full and 0.121 for MIPROv2), so we treat the single grouped split as a headline comparison and the out-of-fold analysis as stability evidence rather than as a deterministic margin claim.
\end{beforebox}

\noindent\textbf{Now}

\begin{nowbox}
The majority baseline has high accuracy because unresolved examples dominate, but zero resolved F1. The handwritten prompt is conservative: it has perfect unresolved precision but misses most gaps. MIPROv2 improves substantially, GEPA-scalar improves recall, and GEPA-category is too coarse in this run. GEPA-full is strongest overall, improving macro-F1 by 10.4 points over MIPROv2 and 8.1 points over GEPA-scalar. A TF-IDF logistic-regression classifier trained on the same raw evidence reaches only 0.523 macro-F1, while an oracle-summary variant reaches 0.858, indicating that lexical supervision does not explain the GEPA-full gain and that evidence interpretation is the bottleneck.\par
\par
We also compute five-fold out-of-fold results over review-thread groups. Mean macro-F1 is 0.698 for GEPA-full, 0.647 for MIPROv2, 0.609 for GEPA-scalar, 0.545 for GEPA-category, and 0.446 for the handwritten baseline. The fold standard deviations are nontrivial (0.151 for GEPA-full and 0.121 for MIPROv2), so we treat the single grouped split as a headline comparison and the out-of-fold analysis as stability evidence rather than as a deterministic margin claim. To test possible feedback leakage, we regenerate full textual feedback with an independent \textbackslash{}texttt\{gpt-5.4-mini\} writer that sees only test-time fields plus the gold label, not the original adjudication rationales; GEPA with this feedback reaches 0.750 macro-F1, nearly matching the original 0.756.
\end{nowbox}

\end{document}
